% !TeX root = ../../Book2.tex
% 纲要标题：### 4.3 有限生成模的结构
% 纲要位置：计划纲要.md，第 930 行。

\section{有限生成模的结构}
\label{b2:sec:0403}

% 纲要要点（供目录核对）：
% * 挠部分与自由部分
% * 不变因子分解与初等因子分解
% * 第一卷交换群定理的统一解释

矩阵的标准形已经给出了有限展示模的分解。要覆盖所有有限生成模，尚需证明有限个生成元之间的关系可以由有限多个关系生成；一般环上这不成立，在主理想整环上则来自自由模的子模定理。分解存在后，还要从模本身恢复分类数据，避免把“展示矩阵不变”误当作“更换展示也不变”。

\subsection{挠部分与自由部分}
\label{b2:subsec:0403:01}

\begin{theorem}[主理想整环上的有限秩子模]\label{b2:thm:pid-free-submodule}
设 \(R\) 为主理想整环，\(N\le R^n\)，其中 \(n\) 有限。则 \(N\) 是自由模，并有一组至多含 \(n\) 个元素的基。
\end{theorem}
\begin{proof}
对 \(n\) 归纳。\(n=0\) 时 \(N=0\)，空集为基。设 \(n>0\)，令 \(\pi:R^n\to R\) 为第一坐标投影。像 \(\pi(N)\) 是 \(R\) 的理想，故写为 \((a)\)。若 \(a=0\)，\(N\) 包含在后 \(n-1\) 个坐标组成的自由模中，应用归纳假设即可。

若 \(a\ne0\)，选 \(v\in N\) 满足 \(\pi(v)=a\)。令 \(N_0=N\cap\ker\pi\)。归纳假设给出 \(N_0\) 的基 \(v_2,\ldots,v_r\)，其中 \(r-1\le n-1\)。任意 \(x\in N\) 的第一坐标为 \(ba\)，故 \(x-bv\in N_0\)，说明 \(v,v_2,\ldots,v_r\) 生成 \(N\)。若
\[
cv+c_2v_2+\cdots+c_rv_r=0,
\]
投影得到 \(ca=0\)。由于 \(a\ne0\) 且 \(R\) 为整环，\(c=0\)，再用 \(N_0\) 中的基独立性得到其余系数全零。因此所列元素是一组基，大小至多为 \(n\)。
\end{proof}

主理想条件在这里把投影像变成一个生成元，整环条件则使该生成元可以消去。证明没有把子模预先当作自由模，也不需要一般模的分类定理。

\begin{theorem}[有限生成模的不变因子分解]\label{b2:thm:pid-module-structure}
设 \(M\) 是主理想整环 \(R\) 上的有限生成模。则存在整数 \(f,s\ge0\) 与非零非单位元素 \(d_1,\ldots,d_s\)，满足 \(d_1\mid\cdots\mid d_s\)，使
\[
M\cong R^f\oplus R/(d_1)\oplus\cdots\oplus R/(d_s).
\]
其中挠子模为循环商分量的直和，而 \(M/t_R(M)\cong R^f\)。
\end{theorem}
\begin{proof}
选一组有限生成元，由\cref{b2:prop:module-presentation}得到满同态 \(q:R^n\to M\)。其核 \(N\) 由\cref{b2:thm:pid-free-submodule}具有有限基，设基数为 \(r\)。把这些基向量写成 \(R^n\) 的列向量，得到矩阵 \(A:R^r\to R^n\)，像恰为 \(N\)。因此 \(M\cong\operatorname{coker}A\)。

对 \(A\) 用\cref{b2:thm:smith-normal-form}与\cref{b2:prop:smith-cokernel}，得到自由坐标与非零对角因子对应的循环商。删去单位因子对应的零分量，便得到题中形式；余下的整除关系仍成立。循环商被非零标量零化，而自由模无挠，所以一个元素是挠元，当且仅当其全部自由坐标为零。这给出挠子模及无挠商的描述。
\end{proof}

这个定理立即说明，主理想整环上的有限生成无挠模是自由模。有限生成的条件不可省略：\(\mathbb Q\) 作为 \(\mathbb Z\) 模无挠，但不自由。若它有非空基，任取基元 \(b\)，\(b/2\) 的有限整数展开乘二后应等于 \(b\)，与基坐标中一个奇数系数等于全体偶数系数冲突；空基只能生成零模。

\ProseParagraphBreak
还须区别分解的类型与补模的选择。挠子模由定义确定，自由补模通常不唯一。在 \(\mathbb Z\oplus\mathbb Z/2\mathbb Z\) 中，\((1,\bar0)\) 与 \((1,\bar1)\) 各自生成一个与挠子模互补的自由子模，两者不同。下面证明自由秩与循环因子的唯一性，这不会使具体补模也变成唯一。

\subsection{不变因子分解与初等因子分解}
\label{b2:subsec:0403:02}

先回顾这里需要的主理想整环算术。每个非零非单位都有不可约因子，否则不断选取非单位真因子，便得到严格增大的主理想升链，与\cref{b2:lem:pid-ascending-chain}矛盾。依次除去不可约因子，若不能在有限步后结束，同样产生这样的升链，所以分解有限。若不可约元 \(p\) 不整除 \(a\)，则 \((p,a)=R\)，由 Bézout 等式可知 \(p\mid ab\) 蕴含 \(p\mid b\)。因此不可约元是素元，素因子分解在相伴与排列意义下唯一。

\ProseParagraphBreak
若 \((a,b)=R\)，则自然同态
\[
R/(ab)\longrightarrow R/(a)\oplus R/(b)
\]
是同构。事实上，取 \(ua+vb=1\)，\(xvb+yua\) 分别同余于 \(x\pmod a\)、\(y\pmod b\)，故映射满射；同时属于 \((a)\) 与 \((b)\) 的元素写成 \(az\)，互素条件使 \(b\mid z\)，因而核正是 \((ab)\)。这给出所需的中国剩余分解。反复使用它，每个循环商可以拆成素元幂对应的循环商。

\begin{definition}\label{b2:def:pid-primary-module}
对不可约元 \(p\)，模 \(M\) 的 \(p\) \term[mo-zhufenliang]{主分量}[primary component of a module]定义为
\[
M_p=\{\,m\in M\mid p^jm=0\text{ 对某个整数 }j\ge1\,\}.
\]
它是子模，因为两个元素分别被 \(p^a,p^b\) 零化时，其和与差被 \(p^{\max(a,b)}\) 零化。循环商 \(R/(p^e)\) 中的 \(p^e\) 称为一个\term[chudeng-yinzi]{初等因子}[elementary divisor]，其中 \(e\ge1\)。
\end{definition}

\begin{theorem}[初等因子及分类唯一性]\label{b2:thm:pid-elementary-divisors}
主理想整环上的有限生成模具有分解
\[
M\cong R^f\oplus\bigoplus_p M_p,
\qquad
M_p\cong\bigoplus_{h=1}^{s_p}R/(p^{e_{p,h}}),
\]
其中只有有限多个 \(M_p\) 非零，指数 \(e_{p,h}\) 为正整数。自由秩 \(f\)、各主分量中的指数多重集，以及\cref{b2:thm:pid-module-structure}中非单位不变因子的相伴类均由 \(M\) 的同构类型唯一确定。
\end{theorem}
\begin{proof}
把\cref{b2:thm:pid-module-structure}中的各个 \(d_i\) 分解为素元幂，再用刚证明的中国剩余分解，得到所需循环直和。\(R^f\) 无挠。在 \(R/(q^e)\) 中，若 \(p\) 与 \(q\) 不相伴，则 \((p,q^e)=R\)，所以乘 \(p\) 是可逆映射，没有非零元素被 \(p\) 的幂零化；在 \(R/(p^e)\) 中则全部元素被 \(p^e\) 零化。因此所得 \(p\) 分量恰是定义中的 \(M_p\)，不依赖分解的选择。

模同构保持挠子模，故也诱导无挠商同构。若两个分解中的自由秩为 \(f,f'\)，得到 \(R^f\cong R^{f'}\)，由\cref{b2:thm:free-rank-invariance}，因为 \(R\ne0\) 且交换，必有 \(f=f'\)。

固定 \(p\)。\(R/(p)\) 是域：任何不属于 \((p)\) 的元素与 \(p\) 互素，Bézout 等式给出其逆元。对 \(j\ge1\)，商模
\[
p^{j-1}M_p/p^jM_p
\]
被 \(p\) 零化，故具有 \(R/(p)\) 向量空间结构。设其维数为 \(c_{p,j}\)。在一个分量 \(R/(p^e)\) 中，若 \(j>e\)，这项层商为零；若 \(j\le e\)，它由 \(p^{j-1}\) 的陪集生成，且系数 \(a\) 使它为零恰好是 \(p\mid a\)，所以同构于 \(R/(p)\)。于是
\[
c_{p,j}=\#\{\,h\mid e_{p,h}\ge j\,\}.
\]
因此指数恰等于 \(j\) 的分量数为 \(c_{p,j}-c_{p,j+1}\)。这些维数只用到 \(M_p\) 及标量作用，所以由模本身确定，恢复全部指数多重集。

最后，将每个 \(p\) 的正指数按非降顺序排列，令 \(s=\max_p s_p\)，把较短列表在左端补零至长 \(s\)。按列相乘 \(d_i=\prod_p p^{e_{p,i}}\)，得到 \(d_1\mid\cdots\mid d_s\)，且每项非单位；中国剩余同构把这些循环商重新组装成相同的主分量。反过来，任何整除链中的每个素元指数都非降，去掉零指数后正是已确定的正指数列表，所以只能用上述左端补零的方式恢复。故不变因子也在相伴意义下唯一。若没有任何非零主分量，取 \(s=0\)，只剩已确定的自由部分。
\end{proof}
\symidx{\(M_p\)：模的 \(p\) 主分量}

这项证明比较的是模自身的无挠商和主分量层商，所以适用于不同生成元数、不同关系数的展示。Smith 对角形中的单位项可以随冗余展示增减，但它们所给的零模不参与这里的非单位因子列表。

\ProseParagraphBreak
例如，若某个主分量的连续层商维数依次为 \(3,2,2,1,0\)，则指数至少为一、二、三、四的分量分别有三、二、二、一个。相邻相减得到指数恰为一、二、三、四的分量数为 \(1,0,1,1\)，故
\[
M_p\cong R/(p)\oplus R/(p^3)\oplus R/(p^4).
\]
这里无需先知道展示矩阵，就能从标量 \(p\) 的连续作用恢复初等因子。

\subsection{交换群结构定理的模论解释}
\label{b2:subsec:0403:03}

取 \(R=\mathbb Z\)，模就是交换群，有限生成模就是有限生成交换群。素元可取正素数，循环商 \(\mathbb Z/(p^e)\) 是阶为 \(p^e\) 的循环群。因此\cref{b2:thm:pid-elementary-divisors}给出了第一卷交换群结构定理的模论解释：整数倍运算使其成为主理想整环上的同一分类问题。

\ProseParagraphBreak
例如，一个有限群的初等因子分解为
\[
\mathbb Z/4\mathbb Z\oplus\mathbb Z/8\mathbb Z
\oplus\mathbb Z/3\mathbb Z\oplus\mathbb Z/9\mathbb Z
\oplus\mathbb Z/5\mathbb Z.
\]
二、三主分量各有两个循环分量，五主分量只有一个。按非降指数排列，并在五主分量前补一个零指数，逐列乘积为 \(4\cdot3=12\) 与 \(8\cdot9\cdot5=360\)，所以不变因子分解为
\[
\mathbb Z/12\mathbb Z\oplus\mathbb Z/360\mathbb Z.
\]
若另有三个无限循环分量，就在两种写法中同时加上 \(\mathbb Z^3\)，有限挠部分的计算不变。

\ProseParagraphBreak
把整数环换为 \(k[t]\)，素数换为首一不可约多项式，整除链与主分量层商的证明仍然适用。区别在于 \(k[t]/(p)\) 的域结构及其 \(k\) 维数：若 \(p\) 的次数大于一，一个初等循环分量的底层 \(k\) 维数便不等于指数本身。下一节将把这种模分类直接用于线性算子，因而同时适用于有限域、实数域及一般基域。
